\section{Evaluation Plan}
\label{sec:evaluation}

This draft does not yet report empirical results. The planned evaluation should
test whether observer-declared timelines improve recovery, auditability, and
agent cooperation in realistic workflows.

\subsection{Case studies}

\begin{itemize}[leftmargin=*]
  \item \textbf{Remote agent work.} An agent performs a task on a server while
        the user later reviews the work from a local GUI.
  \item \textbf{Multi-machine sync.} Facts from local and remote runtimes are
        imported with different freshness boundaries and causal completeness.
  \item \textbf{Replay and debugging.} A user or agent reconstructs why a task
        reached a given state after partial failure.
  \item \textbf{Release evidence.} Build and release facts are shown as a
        timeline with source provenance and verification status.
\end{itemize}

\subsection{Measures}

Potential measures include:

\begin{itemize}[leftmargin=*]
  \item recovery time after interruption;
  \item number of hidden-ordering assumptions discovered during audit;
  \item ability for an agent to reproduce a displayed timeline from exported
        metadata;
  \item correctness of causal-order preservation under clock skew;
  \item rate of degraded-evidence states surfaced rather than silently hidden;
        and
  \item user and agent ability to decide the next safe action from the view.
\end{itemize}

\subsection{Artifacts}

Evaluation artifacts should include fixture ledgers, projected timeline
exports, KFD-4 observer-perspective declarations, and KFD-2 trust assessments.
The paper should not rely on private operational logs for evidence unless they
are sanitized, permissioned, and explicitly marked.
